% PDF 44--55；原书第 13--24 页。
\section{大语言模型实践}
像任何新技术一样，大语言模型在走向成熟乃至普及的过程中，研究与应用阶段有很大的重叠。因而，找到最有价值或最有用的研究文献，或者为某个应用选出最合适的方法，都可能颇为困难。下面几节将梳理大语言模型在实际工作中的不同用法，帮助读者作出判断。

\begin{figure}[!htbp]
\centering
% 重绘自原书第 14 页的图 1.6。
\begin{tikzpicture}[
  font=\footnotesize\sffamily,
  flow/.style={draw=black!50,-{Stealth[length=1.6mm]},line width=.55pt},
  choice/.style={diamond,draw=black!40,aspect=1.54,minimum height=12mm,align=center,inner sep=1.4pt},
  process/.style={draw=black!40,rounded corners=2mm,fill=green!22!gray!18,
    minimum width=18mm,minimum height=9mm,text width=16mm,align=center,inner sep=1pt},
  stage/.style={font=\footnotesize\sffamily\bfseries},
  app/.style={align=center,inner sep=1pt}]
\node[stage] (scope) at (.56,5.75) {确定范围};
\node[stage] (pretrain) at (3.0,5.75) {预训练};
\node[stage] (adapt) at (6.53,5.75) {适配与对齐};
\node[stage] (deploy) at (9.95,5.75) {应用};
\draw[flow,dashed] (scope.east) -- (pretrain.west);
\draw[flow,dashed] (pretrain.east) -- (adapt.west);
\draw[flow,dashed] (adapt.east) -- (deploy.west);

\node[choice,fill=orange!50] (llm) at (.65,2.74) {模型？};
\node[process,fill=yellow!25,text width=19mm] (scratch) at (2.65,4.55)
  {从头预训练\\[-1pt]{\scriptsize 图书、网页语料、代码}};
\node[process,fill=yellow!45,text width=19mm] (ready) at (2.65,.76)
  {预训练大语言模型\\[-1pt]{\scriptsize OpenAI、Llama、Claude}};
\node[choice,fill=green!18!gray!14] (need) at (4.34,2.74) {适配？};
\node[process] (prompt) at (6.47,4.55) {提示工程};
\node[process] (finetune) at (6.47,2.74) {微调};
\node[process] (align) at (6.47,.87) {对齐调优};
\node[choice,fill=blue!16!gray!10,aspect=1.66] (augment) at (8.85,2.74) {增强？};
\node[app] (user) at (10.55,4.55) {查询提示};
\node[process,fill=blue!37!cyan!30,minimum width=9mm,text width=8mm] (rag) at (10.55,3.45) {RAG};
\node[app,draw=blue!55!black,cylinder,shape border rotate=90,
  aspect=.35,minimum height=7mm,minimum width=12mm,fill=blue!12] (db) at (10.55,2.06) {向量库};
\node[app,draw=blue!55!black,minimum height=7mm] (docs) at (10.55,.7) {自定义文档};

\draw[flow] (llm.north) |- (scratch.west);
\draw[flow] (llm.south) |- (ready.west);
\draw[flow] (scratch.east) -| (need.north);
\draw[flow] (ready.east) -| (need.south);
\coordinate (split) at (5.40,2.74);
\draw[draw=black!50,line width=.55pt] (need.east) -- (split);
\draw[flow] (split) |- (prompt.west);
\draw[flow] (split) -- (finetune.west);
\draw[flow] (split) |- (align.west);
\draw[flow] (prompt.south) -- (finetune.north);
\draw[flow] (finetune.south) -- (align.north);
\draw[flow] (prompt.east) -| (augment.north);
\draw[flow] (finetune.east) -- (augment.west);
\draw[flow] (align.east) -| (augment.south);
\draw[flow] (augment.east) -- (10.55,2.74) -- (rag.south);
\draw[flow] (user.south) -- (rag.north);
\draw[flow] (db.north) -- (10.55,2.74);
\draw[flow] (docs.north) -- (db.south);
\end{tikzpicture}

\caption{从应用出发规划大语言模型的开发、适配和使用时，可以考虑的各种路径。}
\label{fig:ch01-paths}
\end{figure}

大语言模型的实践大体涉及三个领域。
\begin{itemize}
\item \textbf{开发（development）}关注如何从头构建大语言模型，包括从数据集、学习目标和网络架构几个角度考虑预训练。
\item \textbf{适配（adaptation）}关注如何修改预训练的大语言模型，使其产生更符合预期的结果。这些结果往往要结合具体应用衡量，应用也会影响选择哪种适配方法。指令调优、对齐和微调是这里的重要主题。
\item \textbf{使用（utilization）}关注如何与大语言模型交互，以发挥其有用的能力。终端用户提示（end-user prompting）和应用开发是这里的重点。
\end{itemize}

本节先概述大语言模型开发，第 2 章再作详细讨论。接着介绍最常用的适配方法，包括对齐调优和参数高效微调，第 4 和第 5 章会展开这些主题。最后，从使用大语言模型的角度考察终端用户提示和应用。后文还会详谈这些内容，相关材料主要见第 6、7、8 章。下面分别介绍每一领域的基本概念。

\subsection{大语言模型开发}
可以从头预训练自己的大语言模型，也可以在已有模型上进行适应性预训练（adaptive pre-training）。由于这些做法复杂，前期成本也高，预计它们不会成为特别常见的开发路径。若团队需要模型掌握专业领域知识，希望更明确地控制偏见、能力或局限，或者希望降低长期的总拥有成本（Total Cost of Ownership，TCO），就可能选择这条路径，而不用已有的众多开源或闭源预训练模型。本节介绍这些做法所涉及的主要因素。

\subsubsection{大语言模型预训练}
大语言模型的许多出色能力来自预训练过程。预训练使模型从大规模数据集中习得基本的语言技能和知识（\refcite{ch01-zhao}{Zhao et al., 2023}）。下面回顾预训练的三个方面，即常用的数据集、网络架构和学习目标。
\begin{itemize}
\item \textbf{数据集。}与开发大语言模型所用的学习目标和网络架构相比，数据集在不同模型间的差异较小。训练通用大语言模型时，通常组合使用网页抓取资料（如 Common Crawl）、维基百科（Wikipedia）、图书语料和代码。代码数据已被证明能显著提高模型的推理能力（\refcite{ch01-ma}{Ma et al., 2023}）。这些来源不同的内容还要经过筛选整理，统一质量标准，并尽量减轻数据重复、有害或仇恨内容等因素的负面影响。确认数据集覆盖范围和质量足够后，开发团队按需要对数据作规范化处理和词元切分（tokenization）。完成这些步骤，数据才可用于预训练。第 2.4 节会更详细地讨论这些内容。

\item \textbf{网络架构。}大语言模型最常用的是 Transformer 架构，它在开创性的论文《注意力就是你所需要的一切》（\textit{Attention Is All You Need}）中提出（\refcite{ch01-vaswani}{Vaswani et al., 2017}）。原始架构由编码器和解码器组成，两者都采用自注意力机制（self-attention mechanism）和前馈神经网络（feed-forward neural network）。第 2.5.3 节将详细介绍研究者怎样改进原始 Transformer 架构。这里选取几个影响较大的模型，说明架构的改进与模型能力之间的关系。
\begin{itemize}[label=--]
\item 纯编码器架构（encoder-only architecture）用于得到输入序列的上下文表示（contextual representation）。这种表示对分类、命名实体识别等多种自然语言处理任务很有用。BERT 是采用这种设计的一个常用模型（\refcite{ch01-devlin}{Devlin et al., 2019}）。
\item 编码器—解码器架构用于机器翻译、问答等序列到序列（sequence-to-sequence）任务。T5 是采用这种架构的一个常用模型（\refcite{ch01-raffel}{Raffel et al., 2020}）。
\item 当学习目标是自回归序列生成时，使用因果解码器架构（causal decoder architecture）。将注意力限制在一个方向上，就可以实现序列生成。采用这种架构的模型擅长文本生成，最熟悉的是 GPT 系列（例如 \refcite{ch01-brown-2020}{Brown et al., 2020}）。
\item 前缀解码器架构（prefix-decoder architecture），又称非因果解码器架构（non-causal decoder architecture），是上述因果解码器的变体。关键差异在于，对输入序列，即前缀，应用双向注意力机制。生成的序列仍采用单向注意力，生成过程仍是自回归的。书中以 Google 的 PaLM 为例，说明这种架构尤其擅长机器翻译等可受益于双向编码的任务（\refcite{ch01-chowdhery}{Chowdhery et al., 2022}）。
\end{itemize}

\item \textbf{学习目标。}预训练所用的学习目标，对大语言模型最终表现出的涌现能力有很大影响。通常，预训练目标旨在尽可能提高自然语言理解和连贯生成的能力。最常用的方法是完整语言建模（full language modeling），即给定此前的词元，以自回归方式预测序列中的下一个词元（\refcite{ch01-zhao}{Zhao et al., 2023}）。其他重要目标包括去噪自编码（denoising autoencoding）与混合去噪器（mixture-of-denoisers）。前者先破坏输入文本序列，再训练网络恢复受破坏的片段（\refcite{ch01-raffel}{Raffel et al., 2020}）。后者组合使用擅长不同任务的去噪器，利用三种主要的预训练任务，分别为标准语言建模、恢复带少量噪声的短片段，以及恢复长片段或高噪声片段（\refcite{ch01-chowdhery}{Chowdhery et al., 2022}）。前者用于常见的序列到序列大语言模型 T5，原书将后者列为 Google 的 PaLM 所用方法。
\end{itemize}

\subsubsection{大语言模型的适应性预训练}
大语言模型通常使用上一节所述的通用数据集预训练，而领域适应性预训练（domain-adaptive pre-training）使用更贴近特定领域的数据集继续训练模型（\refcite{ch01-gururangan}{Gururangan et al., 2020}）。它的目标是使大语言模型的能力更适应有专门词汇或语言用法的领域。例如，生物学有大量术语，物种名称和解剖描述还会使用拉丁语命名法。只在通用数据集上训练的大语言模型，可能对生物学术语掌握得不够充分。使用来自 bioRxiv 之类的数据集适配基础模型，可以克服这些局限。bioRxiv 是面向生物科学文献的专业预印本平台（\refcite{ch01-zhang-a}{Zhang et al., 2024a}）。

\begin{figure}[!htbp]
\centering
% 重绘自原书第 17 页的图 1.7。
\begin{tikzpicture}[font=\small\sffamily]
\begin{scope}
\clip (-3,-1.75) rectangle (3,2.35);
\fill[cyan!30] (-2.23,.76) ellipse [x radius=2.94cm,y radius=1.86cm,rotate=-28];
\fill[yellow!72!orange,opacity=.36] (2.42,-.20) ellipse [x radius=2.88cm,y radius=1.68cm,rotate=3];
\fill[gray!55,opacity=.62] (-.6,.67) ellipse [x radius=1.30cm,y radius=.60cm,rotate=-50];
\fill[black] (-.54,.41) ellipse [x radius=.65cm,y radius=.22cm,rotate=-39];
\end{scope}
\node at (-1.76,-.74) {目标领域};
\node[align=center] at (1.78,-.65) {原大语言模型\\领域};
\node (task) at (.14,2.00) {任务};
\draw[-{Stealth[length=1.9mm]},black!70] (task.south) to[out=-92,in=84] (-.59,.61);
\end{tikzpicture}

\caption{通用大语言模型预训练（原模型领域）、领域适应性预训练（目标领域）和任务微调（任务）涉及的数据分布之间的关系。任务分布所处的浅灰色区域表明，任务通常是较大分布中的一个可观测子集。}
\label{fig:ch01-domain}
\end{figure}

研究还表明，适应性预训练可以显著改善任务微调的效果（\refcite{ch01-gururangan}{Gururangan et al., 2020}），例如特定领域的主题分类，见图~\ref{fig:ch01-domain}。多个领域都观察到了这种效果，金融服务业的一个知名例子是 BloombergGPT。该模型使用混合数据作适应性预训练，其中有新闻报道、财务报告等金融领域数据，也有 Bloomberg 的专有数据。尽管这部分数据只占所用预训练数据总量的 0.7\%，BloombergGPT 在金融领域的多项任务上仍明显优于参数量为它三倍的其他通用大语言模型（\refcite{ch01-wu}{Wu et al., 2023}）。
\transnote{原书第 17 页此处给出 0.7\%，未明确说明分子对应哪一部分数据，也未定义所比较的预训练数据总量。}

具备合适的数据集、足够的算力和充足的预算时，通用预训练与领域适应性预训练各有收益。只要能承担前期成本，以这种方式从头构建大语言模型，就可能在控制输出结果、性能或隐私方面获得显著优势。多数开发团队恐怕很难满足这一条件。下一节将介绍更接近微调的适配方法。这些方法通常只需上述预训练方法所需数据和计算成本的一小部分。

\subsection{大语言模型适配}
大语言模型的涌现能力本身令人瞩目，但即使经过适应性预训练，模型在给定任务上的开箱表现也未必满足每一种应用需求。围绕如何进一步调优大语言模型、使其朝指定结果改进的工具、技术和流程，研究迅速为开发者提供了许多选择。有些方法针对分类或命名实体识别等传统自然语言处理任务进行微调，有些则着重调整生成文本的风格，使之符合开发者基于价值判断形成的审美偏好。无论哪种，和前述预训练方法的主要区别都是使用带标签的训练数据集。本节探讨监督式适配的四个主要方面，包括指令调优、对齐调优、全参数微调和参数高效微调。

\subsubsection{指令调优}
指令调优是一种微调技术，用来增强大语言模型的能力和可控性。它的主要目标是微调通用大语言模型，使之更准确地遵循用户给出的具体指令（\refcite{ch01-zhang-b}{Zhang et al., 2024b}）。做法是将带标签的数据集整理成“（指令，输出）”对，对大语言模型进行监督训练。代码生成、摘要、问答、任务规划或执行等都可以整理成指令调优数据，使开发者能够提高模型在这些任务中遵循指令的能力。

指令调优通常通过全参数微调或参数高效微调来实现。在学习过程中，大语言模型更好地建立起指令与数据对中预期输出之间的联系，生成的回答因此更符合人类指令。不难预料，全参数指令调优的成本可能很高，特别是在用大型指令调优数据集训练超大模型时。因此，LoRA（\refcite{ch01-hu}{Hu et al., 2021}）等参数高效微调方法，以及专门用于指令调优的 HINT（\refcite{ch01-ivison}{Ivison et al., 2023}），都成为可行的方法，可以在所用数据集规模与微调成本之间取得更好的折衷。

\subsubsection{对齐调优}
指令调优旨在让大语言模型遵循人类的具体指令，而对齐调优旨在更广泛地使模型输出符合人类偏好和价值观（\refcite{ch01-lin}{Lin et al., 2023}）。和指令调优一样，原书将对齐调优归为监督微调技术。它高度依赖人工标注者，根据模型回答符合预先确定的偏好或价值观的程度，为回答排序。这涉及标注者本身的主观判断，尤其当他们来自不同的文化或社会背景时。要为微调过程提供一致的对齐反馈，这是一个重要难题，因此必须清楚界定我们希望维护的价值观。对齐调优通常采用以下三个原则。
\begin{itemize}
\item \textbf{有帮助性（helpfulness）}指模型紧扣提示中的指令，帮助用户完成任务的能力。
\item \textbf{诚实性（honesty）}指模型向用户提供准确信息的能力。
\item \textbf{无害性（harmlessness）}指模型避免生成会伤害用户或违背用户价值观与道德观的文本的能力。生成仇恨内容或表现出偏见等，都会降低模型的无害性。
\end{itemize}

第 5 章将详细讨论多种对齐调优方法。其中最知名的两种是基于人类反馈的强化学习（RLHF，见 \refcite{ch01-kaufmann}{Kaufmann et al., 2024}）和直接偏好优化（DPO，见 \refcite{ch01-rafailov}{Rafailov et al., 2023}）。RLHF 使用基于人类判断的反馈拟合反映人类偏好的奖励模型（reward model），再利用该模型微调大语言模型，使奖励最大化。拟合奖励模型和维持对齐训练的稳定性存在困难，DPO 的提出正是为了解决其中一些问题。它直接使用成对的偏好数据，按简单的分类目标（例如最大似然）微调大语言模型。提出者将 DPO 作为一种更简单的对齐调优方法。

\subsubsection{全参数微调}
早期的大语言模型微调主要采用指令调优。研究者将带标签的数据集重新整理成自然语言指令，让大语言模型处理这些指令，更新模型参数（见第 1.5.2.1 节）。然而，用数千条指令样本更新数十亿个参数，需要极高的算力，只有少数拥有足够预算的企业能承担。因此，人们投入很多精力，研究如何提高全参数调优的内存效率和计算效率。提高计算效率最常用的方法是量化（quantization），它在预训练期间或之后减少模型需要的内存（\refcite{ch01-gholami}{Gholami et al., 2021}）。除了用量化提高微调效率，近期研究也表明，降低优化过程的内存消耗既能帮助完成任务，也能减少微调的总成本（\refcite{ch01-lv}{Lv et al., 2023}）。第 4.4 节将进一步讨论这些方法。

\subsubsection{参数高效微调}
另一组大语言模型微调方法试图在提高特定任务表现的同时，尽量减少需要调优的参数数目（\refcite{ch01-zhao}{Zhao et al., 2023}）。下面介绍两种重要的参数高效微调方法。

\textbf{低秩适配（LoRA）}是 \refcite{ch01-hu}{Hu 等人（2021）}提出的方法。它在大语言模型的每个 Transformer 层中注入可训练的秩分解矩阵，同时冻结相应的预训练模型权重，从而减少需要调优的参数数目。注入的矩阵可以提高微调后任务的表现，又不会显著增加推理延迟。下面将看到，其他参数高效微调方法并非都能做到这一点。

\textbf{适配器}是另一种参数高效微调方法。它在每个 Transformer 层的输入和输出部件之间插入小型神经网络模块，以这些模块代替秩分解矩阵。优化适配器的参数时，规模大得多的 Transformer 部件保持固定。适配器显著减少需要调优的参数总量，因而大幅缩短训练时间。不过，在处理流程中加入额外部件会延长推理时间。第 4 章将更全面地讨论这些方法以及其他值得关注的参数高效微调方法。

\subsection{大语言模型的使用}
从最直接的层面看，使用大语言模型主要指终端用户通过提示与模型交互。不过，本节也从应用的层面讨论如何使用大语言模型，即从更高层面利用它。第 8 章会详细介绍这种用法。这里先介绍大语言模型中很有用的一种涌现能力，即上下文学习（\refcite{ch01-wei-2022}{Wei et al., 2022}），再介绍名为思维链提示的进阶方法，说明终端用户编写提示时的一些基本思路。提示工程还有许多创新，本节不逐一讨论，余下各章会介绍细节。从应用角度，本节概述几类主要应用，以及促成这些应用的概念和框架创新。本书后续各章将详谈其中许多主题。第 7 章讨论对话式大语言模型与检索增强生成，第 6 章讨论大语言模型的挑战与评估，第 8 章讨论大语言模型应用的开发与运维。

\subsubsection{上下文学习}
\label{sec:ch01-icl}
为了让大语言模型在任务中发挥涌现能力，往往会借助另一项涌现能力，即上下文学习。这种能力最早在 OpenAI 的早期 GPT 模型上得到展示（\refcite{ch01-brown-2020}{Brown et al., 2020}；\refcite{ch01-radford}{Radford et al., 2019}）。它使大语言模型仅在推理阶段就能从自然语言输入中学习，模型参数无须更新。通常把这类自然语言输入称为提示。根据输入提示中有多少任务示范作为上下文，可以分为零样本提示（zero-shot prompt）、少样本提示（few-shot prompt）或多样本提示（multi-shot prompt）。

\sourcefigure{fig01_08.pdf}{上下文学习演示。对三种参数数量不同的 GPT-3 模型（13 亿、130 亿和 1\,750 亿个参数）分别进行评估，测量它们移除故意插入单词的随机符号时的表现。每个模型都在零样本、单样本和少样本设置下测试。横轴为上下文中的示例数 $k$，纵轴为准确率。虚线表示提示中没有自然语言任务描述时的表现，实线表示把自然语言任务描述作为上下文时的表现。}{fig:ch01-icl}

图~\ref{fig:ch01-icl} 改编自 \refcite{ch01-brown-2020}{Brown 等人（2020）}的图 1.2。它展示了上下文学习对模型完成一项简单任务的影响，即去掉正常单词中插入的随机符号。在被测试的最大 GPT-3 模型（1\,750 亿个参数）上，这种影响最为显著。只给模型一条含自然语言任务描述的提示，就能明显改善表现。若在描述之外仅再给一个任务示范，相对于没有描述的情形，GPT-3 的表现会有更显著的跃升。GPT-3.5、GPT-4、PaLM、Llama-2 等先进模型的这一特性，是持续推动大语言模型应用创新的主要因素之一。怎样设计提示来改善语言模型在指定任务上的表现，第 3 章有详细讨论。原书还将读者引向与终端用户提示相关的第 7 章。
\transnote{原书第 21 页此处指向第 7 章。题为“终端用户提示”的专节是第 4.5 节。}

\subsubsection{思维链推理}
上下文学习的一个值得关注的特例是思维链提示（Chain-of-Thought Prompting，CoT）。这种方法组织提示中的上下文、指令和查询，引导大语言模型按有逻辑的一系列步骤来解决一个较大的任务（\refcite{ch01-wei-2023}{Wei et al., 2023}）。图~\ref{fig:ch01-cot} 演示了思维链的思想。与标准提示相比，在示范中加入逐步推理，可以让大语言模型更好地回答示例问题。上下文学习及其特例思维链提示，都属于在推理阶段用自然语言提示来引出大语言模型的知识、语言理解与推理能力的做法。

\begin{figure}[!htbp]
\centering
% 重绘自原书第 22 页的图 1.9。两个版本都保留例题与两个回答中的数值。
\begingroup
\small
\tcbset{cotpanel/.style={arc=2mm,colback=white,colframe=black!40,
  colbacktitle=black!82,coltitle=white,boxrule=.45pt,
  fonttitle=\sffamily\bfseries\small,left=4pt,right=4pt,top=4pt,bottom=4pt,
  before skip=5pt,after skip=5pt}}
\begin{minipage}[t]{.466\linewidth}
\centering{\sffamily\bfseries 标准提示}\par
\begin{tcolorbox}[cotpanel,title={模型输入}]
\textbf{问}\enspace Bob 举办聚会，买了 3 张披萨，每张切成 12 片。参加聚会的人吃掉了两整张披萨，以及另一张的三分之一。还剩多少片披萨？

\textbf{答}\enspace 答案是 8。

\textbf{问}\enspace Rebecca 举办聚会，买了 4 盒甜甜圈，每盒 12 个。聚会上有 12 人，每人吃 2 个甜甜圈。还剩多少盒甜甜圈？
\end{tcolorbox}
\begin{tcolorbox}[cotpanel,title={模型输出}]
\textbf{答}\enspace 答案是 24。\par
\hfill\textcolor{red!75!black}{\(\times\)}
\end{tcolorbox}
\end{minipage}\hfill
\begin{minipage}[t]{.502\linewidth}
\centering{\sffamily\bfseries 思维链提示}\par
\begin{tcolorbox}[cotpanel,title={模型输入}]
\textbf{问}\enspace Bob 举办聚会，买了 3 张披萨，每张切成 12 片。参加聚会的人吃掉了两整张披萨，以及另一张的三分之一。还剩多少片披萨？

\textbf{答}\enspace Bob 一开始有 3 张披萨，每张 12 片，也就是 \(12+12+12=36\) 片。吃掉两整张，还剩 \(36-12-12=12\) 片。最后一张又吃了三分之一，所以剩 \(12-4=8\) 片，因为 \(12/3=4\)。答案是 8。

\textbf{问}\enspace Rebecca 举办聚会，买了 4 盒甜甜圈，每盒 12 个。聚会上有 12 人，每人吃 2 个甜甜圈。还剩多少盒甜甜圈？
\end{tcolorbox}
\begin{tcolorbox}[cotpanel,title={模型输出}]
\textbf{答}\enspace 4 盒甜甜圈每盒 12 个，共有 \(12+12+12+12=48\) 个。12 人每人吃 2 个，还剩 \(48-24=24\) 个。每盒 12 个，所以还剩 2 盒甜甜圈。\par
\hfill\textcolor{green!50!black}{\bfseries 正确}
\end{tcolorbox}
\end{minipage}
\endgroup

\caption{假想的标准提示和思维链提示的比较。给大语言模型增加推理说明，既能提高答案的准确率，也使模型得出答案的过程更容易解释。}
\label{fig:ch01-cot}
\end{figure}

\subsubsection{应用}
如本章开头所言，公众最熟悉的大语言模型应用，是让模型写作或辅助写作。这类应用通常使用自回归模型。模型从输入文本开始，预测下一个词元，把这个词元追加到输入的末尾，然后以加长后的文本重复这个过程，直至完成生成。第~\ref{sec:intro-opening} 节提到的聊天机器人也以此为基础。人们还用这种方法开发辅助写作软件，帮助起草电子邮件、法律文书、技术手册、数据报告，以及几乎任何可以想到的文字材料。要为某个场景获得符合要求的高质量输出，通常需要用人工审定的数据，例如问答对，微调基础大语言模型，以引导应用给出所需的答案（见第 5 章）。

辅助写作在计算机编程领域也有显著应用。基本思想与自然语言写作相同，但训练数据换成了多种编程语言的大量代码。编程助手可以实时建议自动补全，可以根据伪代码生成与现有代码库风格相容的函数，也可以根据自然语言指令把数据库中的数据填入仪表板。GitHub Copilot 等应用已在程序员中广泛使用，既可简化软件开发中较烦琐的环节，也能增强他们的创造力和解决一般问题的能力。

对话、搜索或检索领域的另一些应用，尝试减轻大语言模型失效的负面影响，例如模型容易产生幻觉，即对查询给出不符合事实的回答。克服这一问题的一种方法是检索增强生成（RAG）。RAG 将一个大语言模型和一个特定主题的知识库配对，以便从中取得上下文。知识库可以是一份汽车使用手册，也可以是某一法律案件的一组文档。RAG 流程接收用户查询，判断查询是否与指定主题有关，然后搜索配对的知识库，提取与用户问题有关的信息。随后将知识库中相关的上下文与原始查询一同交给大语言模型，由它生成回答。这样，RAG 就能利用大语言模型准确回答与所给上下文有关的问题的能力。在推理时提供上下文，模型就无需事先包含那些信息。第 7 章将详细讨论 RAG 系统。

在重视隐私和数据保护的行业中，RAG 尤为常见。这些行业中的领域知识大多超出通用大语言模型的训练数据分布，因此开箱即用的对话应用价值可能有限。直接用受保护的信息，如患者病史，对模型作领域适配也不可取。用于训练大语言模型的信息可能深嵌在模型权重中，日后回答任何一次用户查询时都有泄露的风险。RAG 使受保护信息可以留在大语言模型的训练数据之外，同时仍属于模型可使用的知识范围。

在法律或自然语言处理研究等领域，大语言模型具有当时先进的实体识别与推理能力。过去，数据科学家团队要花几个月开发高质量的实体识别模型。他们采用耗时耗资源的全监督方法，并配合复杂的业务规则系统，才能围绕识别出的相关实体作准确推理并采取行动。如今，基于大语言模型的应用已经能识别这些实体并据此推理，其能力接近律师资格考试成绩位于前 10\% 的人类律师。不过，这类评估能在多大程度上衡量大语言模型相对于人类的能力，仍不清楚（\refcite{ch01-martinez}{Martínez, 2023}）。

本节只举出了众多应用中的一小部分，应用范围还在不断扩大。本书各处还会讨论更多应用，尤其是第 8 章。各领域的从业者继续探索将大语言模型用于自身工作时的收益与局限，与这些技术有关的创新也将不断加快。大语言模型在许多学习、评估和生成基准上表现突出，人们自然希望把这些能力用于解决企业和消费者的问题。研究人员通过微调或上下文学习等技术，让大语言模型适应各种领域，由此产生了许多发挥其新能力的应用。我们希望本书能成为了解大语言模型及其用法的入门教材和参考资料。
